\subsubsection{Pile au lithium}

\setcounter{equation}{0}
\setcounter{Q}{0}

%sujet E3A, Arts et Metiers, Physiqu-Chimie, MP
Les piles au lithium équipent de nombreux appareils électroniques modernes, notamment les téléphones portables et appareils photographiques.

Ce type de pile est constitué d'une borne positive en dioxyde de manganèse \ce{MnO2} et d'une borne négative en lithium ; l'électrolyte est un sel de lithium (\ce{LiPF6}) dissout dans un solvant organique (carbonate de propylène) et concentré en ions \ce{Li} (milieu acide). Les couples électrochimiques concernés sont respectivement \ce{MnO2}/\ce{MnO(OH)} et \ce{Li^+}/\ce{Li}.

\Q \'{E}crire les réactions intervenant à chaque électrode, en précisant leur nature. En déduire la réaction fictive de la pile ainsi que sa f.e.m. théorique initiale. Pourquoi l’électrolyte est-il un solvant organique ?

\solution{
Anode = siège de l'oxydation, cathode = siège de la réduction.
\begin{center}
\begin{tabular}{ccC{9cm}}
$\ominus$ & anode & \ce{Li (s) = Li^+ (org) + e^-} \\
$\oplus$ & cathode & \ce{MnO2 (s) + e^- + H^+ (org) = MnO(OH) (s)} \\
\hline
& & \ce{Li (s) + MnO2 (s) + H^+ (org) = Li^+ (org) + MnO(OH) (s)}
\end{tabular}\end{center}
Cette réaction a une f.e.m. calculable dans les condtions standard : 
\begin{align*}
e^\circ &= E_\oplus^\circ - E_\ominus^\circ \\
&= 1.01 - (-3.04) = 4.05 ~ \si{\volt}
\end{align*}
Les données nous indiquent que le potentiel standard du lithium est très négatif, on a donc un important écart à la référence électrochimique : \ce{H^+ (aq)} / \ce{H2(g)}. On s'attendrait dans l'eau à avoir une réaction spontanée/exergonique entre le lithium et l'eau.
\begin{equation*}
\ce{Li(s) + H^+ (aq) = Li^+ (aq) + \frac{1}{2} H_2(g)}
\end{equation*}}

\Q En approximant que les potentiels rédox en milieu organique sont similaires à ceux en milieu aqueux, exprimer la constante de réaction $K^\circ$ à \SI{298}{\kelvin}. \'{E}valuer-la et commenter. \emph{Aide : Vous utiliserez le fait que la f.e.m. de la pile vaut $0$ lorsqu'elle a terminé de débiter.}

\solution{Lorsque la pile a terminé de débiter, on a : 
\begin{align*}
E_\ominus &= E_\oplus \\
\implies E^\circ (\ce{Li^+ (aq) / Li (s)}) + \frac{RT}{F} \ln \left( \frac{a (\ce{Li^+}) }{a ( \ce{Li} )} \right) &= 
E^\circ (\ce{MnOOH (s) / MnO2 (s)}) + \frac{RT}{F} \ln \left( \frac{a ( \ce{MnO2} ) a ( \ce{H^+} ) }{a ( \ce{MnO(OH)} )} \right) \\
\implies \frac{a (\ce{Li^+}) a ( \ce{MnO(OH)} )}{a ( \ce{Li} ) a ( \ce{MnO2} ) a ( \ce{H^+} )} &= \exp \left( \frac{F}{RT} ( E^\circ (\ce{MnOOH (s) / MnO2 (s)}) - E^\circ( \ce{Li^+ (aq)} / \ce{Li (s)}) \right)
\end{align*}
Or, on sait qu'à l'équilibre thermodynamique (cf BUT1 - chimie des solutions), la constante d'équilibre s'écrit : 
\begin{align*}
K_r^\circ (T) &= Q_{r, eq} \\
K_r^\circ (T) &= \frac{a (\ce{Li^+}) a ( \ce{MnO(OH)} )}{a ( \ce{Li} ) a ( \ce{MnO2} ) a ( \ce{H^+} )}
\end{align*}
On peut alors exprimer la constante d'équilibre en fonction des potentiels des couples mis en jeu : 
\begin{align*}
K_r^\circ (T)= \exp \left( \frac{F}{RT} ( E^\circ (\ce{MnOOH (s) / MnO2 (s)}) - E^\circ( \ce{Li^+ (aq)} / \ce{Li (s)}) \right)
\end{align*}
\`{A} \SI{298}{\kelvin}, on a :
\begin{align*}
K_r^\circ (\SI{298}{\kelvin}) &= 10^{\frac{E^\circ (\ce{MnOOH (s) / MnO2 (s)}) - E^\circ( \ce{Li^+ (aq)} / \ce{Li (s))} }{\alpha}} \\
&= 10^\frac{1.01 - (-3.04)}{0.06} \approx 10^{68} \gg 1
\end{align*}
On peut donc s'attendre à ce que la réaction fictive de la pile soit TOTALE, ce qui se traduit par la disparition du réactif limitant. A priori, ici, il s'agit du lithium métallique. \\
\textbf{\`{A} retenir} La constante d'équilibre d'un système rédox est donnée par : 
\begin{align*}
K^\circ ( T ) =  \exp \frac{ F(E_c^\circ - E_a^\circ)}{RT}
\end{align*}
où les indices $c$ et $a$ indiquent respectivement la cathode et l'anode.
}

\Q Déterminer la quantité de matière de \ce{Li} disponible, ainsi que le nombre $n_e$ de moles d'électrons que peut transférer la pile. En déduire la quantité d'électricité $Q$ (exprimée en $C$ puis en \si{\ampere\hour}) qu'elle peut fournir.

\solution{On commence par modéliser l'avancement de la réaction à l'anode : 
\begin{center}\begin{tabular}{cccccc}
\si{\mole} & \ce{Li(s)} & \ce{ = } & \ce{Li^+ (org)} & \ce{ + } &  \ce{e^-} \\
$t=0$ & $\frac{m_0}{M_0}$ & & 0 & & 0 \\
$t \to \infty$ & $\frac{m_0}{M_0} - \xi_\infty$ & & $\xi_\infty$ & & $\xi_\infty$ \\
\end{tabular}\end{center}
Puisque la réaction est totale, on a alors $\frac{m_0}{M_0} - \xi_\infty = 0 \implies \xi_\infty = \frac{m_0}{M_0}$. Dans ce système, on a 1 mole d'électron échangé par mole de Li consommé, on peut donc écrire : 
\begin{align*}
Q = 1 F \xi_\infty = F \frac{m_0}{M_0} \quad [\si{\coulomb}]
\end{align*}
On sait qu'un courant d'un ampère débité pendant 1 heure est équivalent à \SI{3600}{\coulomb}, on a donc : 
\begin{align*}
Q &= \frac{F}{3600} \frac{m_0}{M_0} \quad [\si{\ampere\hour}] \\
&= 26.8 \frac{m_0}{M_0} \quad [\si{\ampere\hour}]
\end{align*}
\textbf{AN :} $Q = 28 \cdot 10^3 ~ \si{\coulomb} = 7.8 ~ \si{\ampere\hour}$.
}

\Q Exprimer la capacité spécifique massique $C_m$. Positionner la capacité de cette pile au lithium par rapport à des piles pour lesquelles les capacités massiques (en \si{\ampere\hour\per\kilo\gram}) s'élèvent respectivement à $480$ (\ce{Cd}), $500$ (\ce{Zn}) ou $820$ (\ce{Ag}).

\solution{
La capacité massique du pôle $\ominus$ est donnée par :
\begin{align*}
C_m =  \frac{Q}{m_0} \stackrel{AN}{=} 3.9 ~ \si{\ampere\hour\per\gram} = 3.9 \cdot 10^3 ~ \si{\ampere\hour\per\kilo\gram}
\end{align*}
On remarque que du fait que le lithium est un élement très léger, sa capacité est au moins $5$ fois plus importante que celle obtenue avec d'autres électrodes $\ominus$.
}

\Q Calculer l’autonomie, en année, de la pile. Quand est-elle usée ?
\solution{On a vu que le courant $I$ exprime la variation de charge $Q$ au cours du temps :
\begin{align*}
I = \frac{dQ}{dt} = \frac{Q}{\Delta t} \text{ car le courant est constant}
\end{align*}
On trouve donc que $\Delta t = Q / I \stackrel{AN}{\approx} 28 \cdot 10^7 ~ \si{\second} \approx 8.9 ~ \text{ans}$
Cette pile lithium est idéale pour assurer le fonctionnement d'appareils nécessitant de faible courant (\emph{par ex.} dans les appareils photo réflex) et garantit une grande longévité.
}


\subsection*{Données}
\begin{itemize}
\item masse de l'électrode de lithium : $m_0 = 2.0$ \si{\gram} ;
\item masse molaire du lithium : $M_0 = 6.941$ \si{\gram\per\mole} ;
\item courant débité par la pile : $I = 0.1$ \si{\milli\ampere} ;
\item $\alpha = \frac{RT}{F} \ln 10 = 0.06 ~ \si{\volt}$ à \SI{298}{\kelvin} ; 
\item potentiels standard d'oxydoréduction à \SI{298}{\kelvin} : 
\end{itemize}

\begin{center}
\begin{tabular}{ccc}
\toprule
Couple & \ce{Li+ (aq)}/\ce{Li (s)} & \ce{MnO2 (s)}/\ce{MnO(OH) (s)}\\
\midrule
$E^\circ ~ [\si{\volt\per {ESH}}]$ & -3.04 & 1.01 \\
\bottomrule
\end{tabular}
\end{center}

